\section*{v2.8.0 Spectral Admissibility and Resolvent Gatekeeping Core}
\addcontentsline{toc}{section}{v2.8.0 Spectral Admissibility and Resolvent Gatekeeping Core}

\begin{quote}
\textbf{Normalization note (v2.17.1 local review).}
This block is retained as a \emph{historical spectral-admissibility precursor}. Its gatekeeping role is preserved, but the normalized admissibility logic is now carried by the semantic/tensor and grouped-layer unfolding phases.
\end{quote}

The semigroup realization from v2.7.0 still leaves one decisive structural question
open: which generators are admissible before one even starts evolving the system? In
this version the answer is sharpened by adding a spectral gatekeeping layer. The
state evolution is not accepted merely because it can be written as an abstract flow;
it must also possess a resolvent and spectrum compatible with the ledger discipline,
HC confinement, tagged-source accounting, and probe inheritance.

\subsection*{1. Spectral admissibility principle}
Let $\mathfrak X$ be the admissible state space from the previous realization step and
let $\mathcal G$ denote the effective generator of the dynamics, including its
controlled tagged-source coupling. We call $\mathcal G$ \emph{spectrally admissible}
if there exist $\omega_*>0$ and $C_*>0$ such that for every $\lambda$ with
$\Re \lambda\ge -\omega_*$ the resolvent exists on the admissible sector and obeys
\[
\| (\lambda-\mathcal G)^{-1} \|_{\mathfrak X\to\mathfrak X}
\le \frac{C_*}{1+|\lambda|},
\]
after projection to the HC-confined channel space. This estimate expresses that the
generator does not hide uncontrolled growth modes behind formal notation. The ledger
control from earlier versions is therefore promoted into a spectral entrance test.

\paragraph{Principle 2 (spectral gatekeeping).}
A candidate field equation or abstract operator model is structurally admissible only
if its effective generator is spectrally admissible in the above sense and the
corresponding resolvent bounds are compatible with the tagged-source decomposition.
The proof programme may therefore reject an operator before any detailed PDE analysis
if the spectrum already violates ledger closure.

\subsection*{2. Resolvent--ledger compatibility}
The weighted energy $\mathcal E$ from v2.7.0 should now be stable not only along the
flow but already at the resolvent level. Concretely, admissibility requires a bound
of the schematic form
\[
\mathcal E\bigl((\lambda-\mathcal G)^{-1}F\bigr)
\le C_{\mathrm{res}}
\Bigl( \mathcal E(F)
+ \mathrm{Tag}(F)
+ \mathrm{Match}(F)\Bigr)
\]
for every admissible forcing term $F$ and every $\Re\lambda\ge -\omega_*$. Here
$\mathrm{Tag}(F)$ denotes the explicitly source-tagged part and
$\mathrm{Match}(F)$ the mismatch component inherited from probe comparison. This is
the spectral analogue of the semigroup energy estimate: the resolvent may amplify
forcing, but only in a way that remains ledger-readable.

\subsection*{3. HC confinement as spectral projector}
The HC mechanism is now reinterpreted one step more sharply. Earlier versions treated
it as the structural compensator that confines evolution to the septinion closure.
At spectral level this becomes the statement that there exists a projector
$\Pi_{\mathrm{HC}}$ with
\[
(\lambda-\mathcal G)^{-1}
= \Pi_{\mathrm{HC}}(\lambda-\mathcal G)^{-1}\Pi_{\mathrm{HC}}
+ \mathcal R_{\mathrm{tag}}(\lambda),
\]
where $\mathcal R_{\mathrm{tag}}(\lambda)$ is either negligible in the admissible
regime or explicitly source-tagged. Hence uncontrolled extra channels are not merely
suppressed dynamically; they are prevented from entering the admissible resolvent
class except through a tagged remainder. This gives a clean spectral formulation of
HC confinement.

\subsection*{4. Spectral gap and restart robustness}
If the admissible generator satisfies the above bounds together with a positive gap
between the dominant admissible spectrum and the forbidden sector, then restart
stability acquires a stronger meaning. A checkpoint is not only dynamically safe but
spectrally robust: small reconstruction errors remain trapped inside the same
resolvent class and do not move the system into a hidden unstable branch. In this
sense proof memory is upgraded from semigroup continuity to spectral persistence.

\subsection*{5. Probe-sector comparison spectrum}
The two-probe subsystem inherits the same criterion. Let
$\mathcal G_{\mathrm{probe}}$ be the effective probe generator acting on the probe
state space. Then admissibility requires a comparison estimate
\[
\| (\lambda-\mathcal G_{\mathrm{probe}})^{-1} \| 
\le C_{\mathrm{probe,res}}(1+|\lambda|)^{-1},
\qquad \Re\lambda\ge -\omega_{\mathrm{probe}},
\]
with constants controlled by the same ledger hierarchy and mismatch tags. The search
geometry therefore becomes spectrally testable: a probe construction is not accepted
merely because it heuristically mirrors the main evolution, but because its spectrum
remains subordinate to the same admissible comparison regime.

\subsection*{6. Consequence for future field equations and appendix tools}
The field-equation admission checklist is sharpened again. A candidate model should
now be tested in the following order:
\begin{enumerate}[label=(\arabic*)]
  \item identify state space and admissible channel projector;
  \item construct the energy and ledger readout;
  \item verify semigroup dissipativity modulo tagged sources;
  \item prove spectral admissibility and resolvent compatibility; and
  \item verify probe-spectrum inheritance and restart robustness.
\end{enumerate}
This also refines the appendix tools. Strong tools may now additionally be marked as
resolvent estimator, spectral projector certificate, gap witness, or restart
persistence lemma. These labels distinguish genuinely admissibility-producing tools
from merely formal symbolic aids.

\subsection*{7. v\docversion\ readiness criterion}
The present step should be regarded as successful if the following can now be done:
\begin{itemize}
  \item admissible generators can be filtered by a spectral entrance test;
  \item resolvent bounds can be read in the same ledger language as the semigroup flow;
  \item HC confinement is expressible as a spectral projector statement;
  \item restart stability is strengthened to spectral persistence; and
  \item the probe sector inherits a controlled comparison spectrum.
\end{itemize}

This is the gain of v\docversion. The companion monolith now carries not only
closure, channel specialization, transport dynamics, forward invariance, and
field-energy realization, but also a first spectral admissibility layer. Future
field equations are therefore constrained twice: they must generate the correct
ledger-dissipative semigroup and they must belong to the correct resolvent class
before they are even admitted into the structural programme.


